% Generated draft: requires mathematical and source review.
\documentclass[11pt]{article}
\usepackage[margin=25mm]{geometry}
\usepackage{fontspec}
\setmainfont{Latin Modern Roman}
\usepackage{amsmath,amssymb,mathtools,unicode-math}
\setmathfont{Latin Modern Math}
\usepackage{enumitem,needspace,xurl,hyperref}
\setlength{\parindent}{0pt}
\setlength{\parskip}{5pt}
\setlength{\emergencystretch}{3em}
\newcommand{\N}{\mathbb{N}}
\newcommand{\Z}{\mathbb{Z}}
\newcommand{\Q}{\mathbb{Q}}
\newcommand{\R}{\mathbb{R}}
\newcommand{\C}{\mathbb{C}}
\newcommand{\F}{\mathbb{F}}
\newcommand{\A}{\mathbb{A}}
\newcommand{\PP}{\mathbb{P}}
\newcommand{\E}{\mathbb{E}}
\newcommand{\eps}{\varepsilon}
\newcommand{\dd}{\,\mathrm{d}}
\begin{document}
\section*{UnsolvedMath ID: 30004017}
\textbf{Source problem code:} catalogue entry
\subsection{Definitions and mathematical statement}
The supplied record asks for a derived equivalence involving unipotent representations, but does not specify enough data to determine a mathematical proposition uniquely.

Under the usual interpretation, for a positive prime power \(q\), \(\operatorname{GU}_n(q)\) denotes the finite general unitary group of a nondegenerate Hermitian form on \(\F_{q^2}^{\,n}\). The source does not specify the permitted values of \(n\) and \(q\). In ordinary characteristic-zero character theory, an irreducible character is called unipotent if it occurs in a Deligne--Lusztig virtual character induced from the trivial character of a rational maximal torus. This terminology alone does not specify the representation categories intended here: the coefficient field or ring, its characteristic, and the objects and morphisms of those categories are not supplied.

The expression \(\operatorname{GL}_n(-q)\) has no ordinary finite-field interpretation for positive \(q\), since a finite field cannot have negative cardinality. It may indicate a formal substitution in formulas, but the record does not define that substitution or an associated representation category. No replacement group or category can therefore be specified from the supplied evidence.

For clarity, if partitions of \(n\) are intended as labels, a partition is a nonincreasing sequence of nonnegative integers with finite sum \(n\). The dominance relation is
\[
\lambda\unrhd\mu
\quad\Longleftrightarrow\quad
\sum_{i=1}^{r}\lambda_i\geq\sum_{i=1}^{r}\mu_i
\quad\text{for every }r\geq1,
\]
where trailing parts are zero. The source does not specify the relevant labeling or how an equivalence must respect this order. It also does not define the relevant Lusztig \(a\)-function or identify the Deligne--Lusztig varieties and their geometric realizations. The original wording uses an ``A-function,'' whereas the cleaned statement uses Lusztig's \(a\)-function; their intended relationship is unresolved.

The supplied question can consequently be retained only in the following schematic form. Once the missing categories \(\mathcal C_U\) and \(\mathcal C_L\), their derived-category conventions, and the compatibility conditions have been specified, does there exist a triangulated equivalence
\[
\Phi:\mathcal D(\mathcal C_U)\xrightarrow{\sim}\mathcal D(\mathcal C_L)
\]
that is compatible with the intended dominance order, the intended function, and the indicated Deligne--Lusztig geometric realizations on both sides? Here \(\mathcal D\) is a placeholder for the unspecified derived-category construction, including its boundedness convention. No precise compatibility equations are supplied.

There is a further unresolved source qualification: the supplied literature assessment describes the extracted material as a fragment concerning a derived self-equivalence and Macdonald-polynomial operators, rather than a complete formal question. It supplies neither those operators nor a relation identifying that fragment with the displayed two-category question. This draft therefore records the proposed question and its missing data without asserting that it is a well-posed conjecture.
\subsection{Short English statement}
The record asks whether the stated unitary and linear unipotent representation categories admit a derived equivalence respecting dominance, a specified function, and Deligne--Lusztig geometry. The notation \(\operatorname{GL}_n(-q)\), the categories, and the compatibility requirements are undefined. The supplied background also describes a different, incomplete fragment. A precise conjecture requires these issues to be resolved.
\subsection{Sources}
\begin{itemize}
\item Local UnsolvedMath registry, supplied complete source record, permanent ID \(30004017\), source problem code \texttt{catalogue entry}, title ``Derived Equivalence of Unipotent Unitary and Linear Representations.'' The record gives proposed year \(2018\) and names no proposer.
\item Supplied report citation: \emph{Enveloping Algebras and Geometric Representation Theory} (2019), Oberwolfach Reports, DOI \texttt{10.4171/owr/2018/52}, pp.\ 3081--3146. \url{https://doi.org/10.4171/owr/2018/52}. The background attributes the material to workshop open-problem discussions; no specific page, speaker, or passage locating this question is supplied.
\item Supplied literature source: ``Oberwolfach Report 52/2018.'' \url{https://ems.press/content/serial-article-files/46773}.
\item Literature assessment dated exactly \texttt{2026-08-21}: the supplied curation labels the record open, explaining that the extracted material is an incomplete fragment about a derived self-equivalence and Macdonald-polynomial operators and that its unresolved status was normalized to open. This is supplied dated curation, not an independently verified literature assessment. Its time-sensitive status and cited claims require independent verification.
\item The registry labels the statement \texttt{corrected\_verified} and says notation, typography, duplicate clauses, or omitted definitions were repaired while preserving mathematical content. This is a supplied editorial assertion; the formulation and source qualifications remain unresolved in this draft. Neither linked source was independently consulted for this draft.
\end{itemize}
\end{document}
